% GENERATED FILE. Edit canonical lessons, then run npm run generate:books.
% canonical-source-hash: fnv1a64:0452bbebbe7e6244
% canonical-lessons: MW-C175-panihari, MW-C175-idani, MW-C175-dhovno, MW-C175-kalsiyo, MW-C175-jhari

\chapter{Water for the House}
\label{ch:mw-water-for-the-house}
\hlchaptermodality{175}

\begin{chapteropening}
\textbf{By the end of this chapter:} \emph{I can say who fetches the water and what it is carried in.}

Complete the last lesson of chapter 175: I can say who fetches the water and what it is carried in.
\end{chapteropening}

\section[\texorpdfstring{\mw{पणिहारी} (paṇihārī)}{पणिहारी (paṇihārī)}]{\mw{पणिहारी} (paṇihārī) — a water-carrier; a woman who fetches water}

\label{lesson:MW-C175-panihari}



\begin{quote}
Before the new one: say the Marwadi for to take off, then the Marwadi for a plait.
\end{quote}

\subsection*{You'll want to know: \mw{पणिहारी}}
\textbf{\mw{पणिहारी}} — \emph{paṇihārī} — \textquotedblleft{}a water-carrier; a woman who fetches water\textquotedblright{}.

\begin{quote}
\mw{पणिहारी} \mw{बावड़ी} \mw{सूं} \mw{पाणी} \mw{लावै} \mw{है।}
\end{quote}

\begin{quote}
\emph{paṇihārī bāvṛī sū̃ pāṇī lāvai hai.}
\end{quote}

\begin{quote}
The water-carrier brings water from the stepwell.
\end{quote}

A feminine noun.

\begin{cousinweb}
\textbf{\mw{पणिहारी}} (\emph{paṇihārī}) is built on a reconstructed \emph{*pānīyahāra}, \textquotedblleft{}a water-carrier\textquotedblright{}, from the word behind \textbf{\mw{पाणी}} and a word for carrying; CDIAL lists Old Marwari \emph{pāṇiyārī} (CDIAL 8088).
\end{cousinweb}

\begin{grammarlens}[title={What you've built}]
One more word for saying what you do.
\end{grammarlens}

\subsection*{Your turn}
\begin{itemize}
\raggedright
  \item \emph{Say it:} \emph{paṇihārī}
  \item \emph{Say it:} \emph{paṇihārī}, once more
  \item \emph{Say it:} \emph{paṇihārī bāvṛī sū̃ pāṇī lāvai hai.} — The water-carrier brings water from the stepwell.
  \item \emph{Recall:} say the Marwadi for to take off, then the Marwadi for a plait, then say \emph{paṇihārī} again
\end{itemize}

\begin{tcolorbox}[breakable,colback=teal!4,colframe=teal!35!black,title={Before you move on}]
What does \mw{पणिहारी} mean? (\textquotedblleft{}A water-carrier; a woman who fetches water\textquotedblright{}.) Say it once more.
\end{tcolorbox}
\section[\texorpdfstring{\mw{ईंडाणी} (ī̃ḍāṇī)}{ईंडाणी (ī̃ḍāṇī)}]{\mw{ईंडाणी} (ī̃ḍāṇī) — a head-pad (for carrying pots)}

\label{lesson:MW-C175-idani}



\begin{quote}
Before the new one: say the Marwadi for a plait, then the Marwadi for a water-carrier.
\end{quote}

\subsection*{You'll want to know: \mw{ईंडाणी}}
\textbf{\mw{ईंडाणी}} — \emph{ī̃ḍāṇī} — \textquotedblleft{}a head-pad (for carrying pots)\textquotedblright{}.

\begin{quote}
\mw{ईंडाणी} \mw{नरम} \mw{है।}
\end{quote}

\begin{quote}
\emph{ī̃ḍāṇī naram hai.}
\end{quote}

\begin{quote}
The head-pad is soft.
\end{quote}

A feminine noun.

\begin{cousinweb}
\textbf{\mw{ईंडाणी}} (\emph{ī̃ḍāṇī}) is inherited from Sanskrit \emph{iṇḍva}, two grass pads that kept the hands from the fire; CDIAL lists Hindi \emph{ī̃ḍwā}, a ring of rolled grass or cloth that cushions a load on the head (CDIAL 1559). CDIAL gives no Marwari form in this entry.
\end{cousinweb}

\begin{grammarlens}[title={What you've built}]
One more word for saying what you do.
\end{grammarlens}

\subsection*{Your turn}
\begin{itemize}
\raggedright
  \item \emph{Say it:} \emph{ī̃ḍāṇī}
  \item \emph{Say it:} \emph{ī̃ḍāṇī}, once more
  \item \emph{Say it:} \emph{ī̃ḍāṇī naram hai.} — The head-pad is soft.
  \item \emph{Recall:} say the Marwadi for a plait, then the Marwadi for a water-carrier, then say \emph{ī̃ḍāṇī} again
\end{itemize}

\begin{tcolorbox}[breakable,colback=teal!4,colframe=teal!35!black,title={Before you move on}]
What does \mw{ईंडाणी} mean? (\textquotedblleft{}A head-pad (for carrying pots)\textquotedblright{}.) Say it once more.
\end{tcolorbox}
\section[\texorpdfstring{\mw{ढोवणो} (ḍhovṇo)}{ढोवणो (ḍhovṇo)}]{\mw{ढोवणो} (ḍhovṇo) — to carry (a load)}

\label{lesson:MW-C175-dhovno}



\begin{quote}
Before the new one: say the Marwadi for a water-carrier, then the Marwadi for a head-pad.
\end{quote}

\subsection*{You'll want to know: \mw{ढोवणो}}
\textbf{\mw{ढोवणो}} — \emph{ḍhovṇo} — \textquotedblleft{}to carry (a load)\textquotedblright{}.

\begin{quote}
\mw{मां} \mw{घड़ो} \mw{ढोवै} \mw{है।}
\end{quote}

\begin{quote}
\emph{mā̃ ghaṛo ḍhovai hai.}
\end{quote}

\begin{quote}
Mother carries the water pot.
\end{quote}

Past (I …): \mw{म्हैं} \mw{पाणी} \mw{ढोयो}, \emph{mhaĩ pāṇī ḍhoyo}.

\begin{cousinweb}
\textbf{\mw{ढोवणो}} (\emph{ḍhovṇo}) is inherited from Sanskrit \emph{ḍhaukayati}, \textquotedblleft{}brings near\textquotedblright{}; CDIAL lists Hindi \emph{ḍhonā}, \textquotedblleft{}to convey\textquotedblright{} (CDIAL 5610). CDIAL gives no Marwari form in this entry.
\end{cousinweb}

\begin{grammarlens}[title={What you've built}]
One more word for saying what you do.
\end{grammarlens}

\subsection*{Your turn}
\begin{itemize}
\raggedright
  \item \emph{Say it:} \emph{ḍhovṇo}
  \item \emph{Say it:} \emph{ḍhovṇo}, once more
  \item \emph{Say it:} \emph{mā̃ ghaṛo ḍhovai hai.} — Mother carries the water pot.
  \item \emph{Recall:} say the Marwadi for a water-carrier, then the Marwadi for a head-pad, then say \emph{ḍhovṇo} again
\end{itemize}

\begin{tcolorbox}[breakable,colback=teal!4,colframe=teal!35!black,title={Before you move on}]
What does \mw{ढोवणो} mean? (\textquotedblleft{}To carry (a load)\textquotedblright{}.) Say it once more.
\end{tcolorbox}
\section[\texorpdfstring{\mw{कळसियो} (kaḷsiyo)}{कळसियो (kaḷsiyo)}]{\mw{कळसियो} (kaḷsiyo) — a small water pot}

\label{lesson:MW-C175-kalsiyo}



\begin{quote}
Before the new one: say the Marwadi for a head-pad, then the Marwadi for to carry.
\end{quote}

\subsection*{You'll want to know: \mw{कळसियो}}
\textbf{\mw{कळसियो}} — \emph{kaḷsiyo} — \textquotedblleft{}a small water pot\textquotedblright{}.

\begin{quote}
\mw{कळसियो} \mw{मेज} \mw{माथै} \mw{है।}
\end{quote}

\begin{quote}
\emph{kaḷsiyo mej māthai hai.}
\end{quote}

\begin{quote}
The small water pot is on the table.
\end{quote}

A masculine noun.

\begin{cousinweb}
\textbf{\mw{कळसियो}} (\emph{kaḷsiyo}) is inherited from Sanskrit \emph{kalaśa}, \textquotedblleft{}a waterpot\textquotedblright{}; CDIAL lists Old Marwari \emph{kalasa}, \textquotedblleft{}a pitcher\textquotedblright{}, and a Gujarati word for a jug that is close to this one (CDIAL 2920).
\end{cousinweb}

\begin{grammarlens}[title={What you've built}]
One more word for saying what you do.
\end{grammarlens}

\subsection*{Your turn}
\begin{itemize}
\raggedright
  \item \emph{Say it:} \emph{kaḷsiyo}
  \item \emph{Say it:} \emph{kaḷsiyo}, once more
  \item \emph{Say it:} \emph{kaḷsiyo mej māthai hai.} — The small water pot is on the table.
  \item \emph{Recall:} say the Marwadi for a head-pad, then the Marwadi for to carry, then say \emph{kaḷsiyo} again
\end{itemize}

\begin{tcolorbox}[breakable,colback=teal!4,colframe=teal!35!black,title={Before you move on}]
What does \mw{कळसियो} mean? (\textquotedblleft{}A small water pot\textquotedblright{}.) Say it once more.
\end{tcolorbox}
\section[\texorpdfstring{\mw{झारी} (jhārī)}{झारी (jhārī)}]{\mw{झारी} (jhārī) — a spouted water pot}

\label{lesson:MW-C175-jhari}



\begin{quote}
Before the new one: say the Marwadi for to carry, then the Marwadi for a small water pot.
\end{quote}

\subsection*{You'll want to know: \mw{झारी}}
\textbf{\mw{झारी}} — \emph{jhārī} — \textquotedblleft{}a spouted water pot\textquotedblright{}.

\begin{quote}
\mw{झारी} \mw{मांय} \mw{पाणी} \mw{है।}
\end{quote}

\begin{quote}
\emph{jhārī mā̃y pāṇī hai.}
\end{quote}

\begin{quote}
There's water in the spouted pot.
\end{quote}

A feminine noun.

\begin{cousinweb}
\textbf{\mw{झारी}} (\emph{jhārī}) comes from a reconstructed \emph{*jhārikā}, \textquotedblleft{}a water-pot\textquotedblright{}; CDIAL lists Hindi and Gujarati \emph{jhārī}, a spouted pot to drink from (CDIAL 5377). CDIAL gives no Marwari form in this entry.
\end{cousinweb}

\begin{grammarlens}[title={What you've built}]
One more word for saying what you do.
\end{grammarlens}

\subsection*{Your turn}
\begin{itemize}
\raggedright
  \item \emph{Say it:} \emph{jhārī}
  \item \emph{Say it:} \emph{jhārī}, once more
  \item \emph{Say it:} \emph{jhārī mā̃y pāṇī hai.} — There's water in the spouted pot.
  \item \emph{Recall:} say the Marwadi for to carry, then the Marwadi for a small water pot, then say \emph{jhārī} again
\end{itemize}

\begin{tcolorbox}[breakable,colback=teal!4,colframe=teal!35!black,title={Before you move on}]
What does \mw{झारी} mean? (\textquotedblleft{}A spouted water pot\textquotedblright{}.) Say it once more.
\end{tcolorbox}
